\chapter{R8: a joint divergent-$D$ certificate, a paired-$D$ thin lift, and the remote join}
\label{component:pe07041e033da}
\noindent\textbf{Component }\path{pe07041e033da}\par
{\small\noindent Source: \path{R17/proofs/weak_scale/joint_D_R8_closure.tex}\par}
{\small\noindent SHA256: \path{e07041e033da07344e496a0c5af3070bc7e6263375a7929c9b1aa32bb43b806b}\par}
{\small\noindent Inventory use: JD1-JD24 and projective join\par}
\medskip
\section{R8: a joint divergent-$D$ certificate, a paired-$D$ thin lift, and the remote join}
\label{pe07041e033da:sec:r17-R8-joint-D}

\paragraph{Scope and conventions.}
All identities use $F/4$ and the unchanged physical C2 parent, with
$\lambda_{\rm event}=\beta=0$, $d=1$, and $n\ge0$.  Every original
stationary-contact, nonnegative-coefficient, strict-gap, total-scale,
and prior-admission-complement intersection is retained.  The two real
closed source/testing fields are independent.  The symbol $\lambda$
below is the positive circuit coefficient, not an event coefficient.

R8 means $B,z,y,h,j>0$ infinitesimal, $t=y/B$ and $g=G/B$ positive
bounded units, $h=H/B$, $a=z/B$, $r\in(0,1)$, and $e=1-r$.  Put
$x=z=aB$, $l=1-h$.  The exact gap is
\[
                 j=1+xr-l/t.                         \tag{JD1}
\]
In particular $t_0=1$.  This section proves:
\begin{enumerate}
\item every independent test and state on the whole R8 subdomain
      where $r$ is a positive bounded unit;
\item the entire remote testing complement of both angles
      infinitesimal on all of R8, including $r$ infinitesimal.
\end{enumerate}
The simultaneous $r,h,j$-infinitesimal local domain is proved by the
separate joint-$n,A,D$ lemma in
\path{weak_r8_small_r/r8_small_r_local.tex}; it is not
inferred from the present $r$-unit proof.  Its union with the two
assertions above closes the whole ordered R8 shape.  No R1 or R2
admission is asserted here.

The inherited exact inputs are R15 RC1--RC13 and its source appendix,
R16 R10.9--R10.12, TL1--TL10 and the two complete-support checks, the
CA/CS/SE bounded contact conductors, and the critical-wall polynomial
identities SEO10--SEO20.  They are used as explicit identities with
the bounds verified below.  In particular no physical C1 theorem is
invoked at a negative formal beta.

\paragraph{A single contact-zero row removes every naked $D$ pole.}
Use the homogeneous R15/R16 roots
\[
 E=tuW-vZ,\quad N_y=(Y-huW)/\sqrt h,\quad \xi=X+ruZ,
 \quad R_D=v(u^2-x)Z,\quad F_v=(v^2-y)Z.
\]
The other roots $Q,N,T,R,J$ and source coefficients are exactly those
of R10.9.  Define
\[
 L_D={l\over t}E-rR_D-\sqrt hN_y+uv\xi,
 \qquad M_D=\sqrt D\,L_D.                            \tag{JD2}
\]
The original event row satisfies the exact homogeneous identity
\[
 uW+uvX-Y-vZ=L_D-jvZ.                               \tag{JD3}
\]
Thus $L_D$ vanishes at the actual source contact, including its full
state normal.  This is a legal separated root at every testing axis.

Let $\mathscr P$ be the right side of R10.9 with the two exact
positive squares $N^2,T^2$ omitted, and let $\mathscr P_0$ denote
its formal coefficient specialization $D=0$.  The other source
letters are held fixed in this notation; it does not change the
physical source.  Direct collection of the complete $D$ channel gives
\[
 \boxed{\quad \mathscr P=D L_D^2+\mathscr P_0
                        -2Dj\,uF_v\xi.\quad}        \tag{JD4}
\]
Here is the full coefficient calculation.  The hard contribution is
$D(lE/t-\sqrt hN_y+uv\xi)^2$; the remaining $D$ contributions are
$Dr^2R_D^2-2Dr(lE/t-\sqrt hN_y+uv\xi)R_D$ and
$2D\{l/t-1-rx\}uF_v\xi$.  Their sum is exactly JD4 by JD1.
This accounts for every hard and first-normal term, rather than just
the angular graph or its Hessian.

Use the polynomial inverse
\[
 E={t\over l}\{D^{-1/2}M_D+rR_D+\sqrt hN_y-uv\xi\}. \tag{JD5}
\]
On R8, $t/l$ is a unit, $D^{-1/2}$ is infinitesimal, and $Dj$ is
bounded.  Every coefficient of the new remainder is bounded whenever
the remaining source arrays of R10.9 are bounded.  In particular no
ratio $h/j$, $x/j$, or testing angle divided by a source opening occurs.
The full divergent rank-one block has been normalized jointly.

\paragraph{The exact residue after the joint normalization.}
Put $\sigma=t\lambda x/\sqrt q$ and $\mu=m-a^2q$ as in R10.9.
At the R8 source/testing corner, JD4--JD5 have the complete residue
\[
 \begin{split}
 M_D^2+r_0^2(C_0+n_0)R_D^2+(Ah)_0N_y^2
  +(n_0+\tfrac12-\mu_0)\xi^2
  +(c_F)_0F_v^2+\lambda_0J^2+R^2+Q^2-2\sigma_0Q\xi.
 \end{split}                                                   \tag{JD6}
\]
This formula applies when $r,n$ are bounded and $r$ a unit.  Its
derivation can retain $n,C,Ah,\mu,c_F,\lambda,\sigma$ as free source
letters: specialize $u=v=\sqrt h=D^{-1/2}=x=\alpha=0$ and $t=1$
in JD4--JD5.  In particular the $E,R_D,N_y$ terms become
$r^2(C+n)R_D^2+AhN_y^2$, with no missing $D$ contribution.
The separate $N^2,T^2$ remain intact.  The focused checker reconstructs
the original hard polynomial and verifies both JD4 and this full
residue.

\paragraph{Bounded $a$, unit $r$, and unit $e$.}
Positivity of $g=1-ae+hx$ makes $a$ bounded whenever $e$ is a unit.
The unchanged source formulas, with $B=h=0$ and $t=1$ only after
clearing unit denominators, give
\[
 \begin{gathered}
 n_0={e_0\over2r_0(1+a_0)},\quad
 (m/a^2)_0={e_0\over2g_0(1+a_0)},\\
 C_0={1+a_0(2r_0-1)\over2(1+a_0)},\quad
 (Ah)_0=(Dj)_0={1+a_0r_0\over2(1+a_0)},\\
 \lambda_0={r_0(1+a_0r_0)\over2(1+a_0)},\quad
 (c_F)_0={(1+a_0r_0)g_0\over2(1+a_0)}.                
 \end{gathered}
\tag{JD7}
\]
All displayed pivots are positive units.  Indeed
$1+a_0(2r_0-1)=g_0+a_0r_0>0$.  The source cancellation defining
$\lambda=(m-a^2q)/(ta^2B)$ is regular also at $a=0$; no division
by an infinitesimal $a$ survives its canceled formula.
Here $q$ is a unit, $\mu=ta^2B\lambda$ is infinitesimal and
$\sigma$ is infinitesimal.  Therefore JD6 is strictly positive.
All omitted entries have zero residue by the finite polynomial
JD4--JD5 and bounded source arrays.  The bounded separated-Gram
lemma gives the complete local parent for every independent ratio
$u/v$, every state, and both literal axes.

\paragraph{Bounded $a$, infinitesimal $e$, including $n=0$.}
Now $r_0=1$, $C_0=(Ah)_0=1/2$, $n,m,\mu$ are infinitesimal,
$\lambda_0=(c_F)_0=1/2$, and
\[
                    (\alpha/B^2)_0=(1+a_0)/2.         \tag{JD8}
\]
The physical inequality OP2 supplies the crucial source lower bound,
without dividing by $n$:
\[
 {e\over B}\ge
 {(1+x)(h+t)(1+ar)\over P_n/B},\qquad P_n/B\sim1.
                                                               \tag{JD9}
\]
It remains valid when $e/B$ is unbounded and at literal $n=0$.
Since $q\ge\alpha+e^2/2$, the $Q,\xi$ Schur pivot satisfies
\[
 \begin{split}
 \operatorname{st}\left(\tfrac12-\mu
                 -{t^2\lambda^2a^2B^2\over q}\right)
 &\ge {1\over2}-{a_0^2\over2(1+a_0)(2+a_0)}\\
 &= {2+3a_0\over2(1+a_0)(2+a_0)}>0.                  
 \end{split}\tag{JD10}
\]
This is a bound entirely in the source field.  The entry $\sigma$
is bounded by JD8, so the whole block is bounded and strict.
The remaining entries of JD6 are positive units.  This proves all
local tests and all states on this source domain, including $n=0$.

\paragraph{Unbounded $a$: the thick layer.}
Put $s=1/a$, $w=ae$, $b=a^2B$.  The exact R10 source identity
still reads
\[
 n=s^2{wQ_*-bK_*\over2r\Delta_*},\qquad
 Q_*\sim1,\ K_*\sim t,\ \Delta_*\sim1.             \tag{JD11}
\]
Thus $b/w$ is bounded, $b$ is bounded, $g_0=1-w_0>0$, and
\[
\begin{gathered}
 (m/w)_0={1\over2(1-w_0)},\quad C_0=(Ah)_0=1/2,
 \quad (Dj)_0=1/2,\\ \lambda_0=1/2,
 \quad(c_F)_0=(1-w_0+b_0)/2>0.
\end{gathered}                       \tag{JD12}
\]
The source formulas for these quantities have no $j$ denominator;
only $D$ itself is unbounded, and JD4 has removed that pole.
For $\rho=n/e^2$ positive noninfinitesimal, retain the exact R10.10
identities and R10.11 Schur pivot:
\[
 \bar q=a^2q=w^2\{\rho r^2+H\}+a^2\alpha,
 \quad H=m+\tfrac12,\quad \chi=m/w,
 \quad \chi-H\in\operatorname{Inf},
\]
\[
 S=\tfrac12-\mu-{(t\lambda x)^2\over q}
 =H-{m^2\over\bar q}
 \ge {H\rho r^2\over\rho r^2+H}
       +{H^2-\chi^2\over\rho r^2+H}.                 \tag{JD13}
\]
The first term has positive residue and the second is infinitesimal,
also for unbounded $\rho$.  Moreover $x/\sqrt q=b/\sqrt{\bar q}$
is bounded by JD11 and $\bar q\ge Hw^2$.  JD6 consequently has
strict residue on $(M_D,R_D,N_y,\xi,F_v,J,R,Q)$, with the only
possible nonzero cross residue in its strictly positive $Q,\xi$
block.  Restore $N^2,T^2$.  This proves the full local parent, not
just its contact-state restriction.

\paragraph{Unbounded $a$: the thin layer and the exact paired-$D$ repair.}
Assume $0<\rho=n/e^2$ infinitesimal.  JD11 gives
$b/w\sim1/t$, so $\psi=w/b$ is a unit.  The formal C1 contact is
exactly OP4, with native parameters
\[
 x_1=x,\quad\eta_1=t,\quad\zeta_1=s,\quad r_1=h,
 \quad\theta_1=e,\quad t_1=w.
\]
TL2 still gives $\beta_1=-s_n/2$ infinitesimal and
$\kappa=(1+2n)/(1-s_n)\sim1$.  The native gaps satisfy
$M_1=g$ a unit and $J_1=j>0$.  The canceled CA/CS/SE remainders
have no inverse $J_1$; their divergent paired-$D$ mass is retained
in SD3.  The formal source has $C_1,A_1h$ positive units,
$D_1j$ a positive unit, positive $m_1,\alpha_1$, and positive
$\beta_1+1/2$.  These facts follow directly from TL3 and JD12.
Its bounded opening ratio $t_1/e_{1,\mathrm{scale}}^2=w/(tb)$
is a unit.  Thus the explicit CA/CS/SE strict base and their finite
conductors apply exactly as in the R16 thin proof; the reservation
$(2+4\beta_1)\nu-\nu$ is positive because $\beta_1$ is
infinitesimal.  This is an algebraic use of those formulas, not an
application of a physical C1 theorem with negative beta.

Here is the one change needed to the R16 perturbation.  In native
notation put
\[
 T_1=c_a=1+xr,\quad a_1=l/t,\quad
 b_D={J_1\over T_1}={j\over c_a},\quad
 h_1={a_1\over T_1}=1-b_D,
\]
\[
 q_B=v(Z-uX),\quad M_B=(uW-Y)-h_1q_B,
 \quad R_{D,\mathrm{event}}=M_B-b_Dq_B.               \tag{JD14}
\]
The old retained paired-$D$ packet in $\kappa F_1/4$ is
$\kappa(D_1+1/2)h_1M_B^2$.  In the actual parent retain instead
the positive contact-zero packet $(D_2+1/2)h_1M_B^2$.
For $\delta_D=(D_2+1/2)-\kappa(D_1+1/2)$ the entire resulting
difference is
\[
 \boxed{\quad
 \delta_D\{R_{D,\mathrm{event}}^2-h_1M_B^2\}
 =\eta_D\{M_B^2-2M_Bq_B+b_Dq_B^2\},\quad
 \eta_D=\delta_D b_D.\quad}                          \tag{JD15}
\]
By the exact TL3 coefficients, with $\Theta=n/B$,
\[
 \eta_D={\Theta\{S_cP_D/t-jP_\kappa\}
                    \over K(h+t)c_a},\qquad
       \eta_D/\rho=O(w).                            \tag{JD16}
\]
Every denominator is a unit.  In particular the previously unbounded
$\rho w/j$ correction is absent; no relative $w/j$ bound is needed.

Keep the R16 positive contact-zero $n$ row TL4, the actual centered
$A$ packet, and every other unchanged retained packet.  The exact
full stress is TL6 with its $\delta_D R_{D,\mathrm{event}}^2$
line replaced by the right side of JD15.  Both parents and all
retained packets have zero contact value and full first jet, so this
new entire stress has the same exact vanishing.

For completeness its new term satisfies every finite-support premise
of the inherited conductor.  The exact inverse state formulas give
\[
 q_B=v\{D_u-uN-puD_uv\mathfrak w\},\qquad p=t_1/R_1.
                                                               \tag{JD17}
\]
At zero normals $q_B=vD_uV_0$.  Hence the new scalar term is in
$(v)^2$, both complete first normals are in $(v)$, and the remaining
hard entries are bounded.  There are no terms above normal degree
two.  Every dependence on $\mathfrak w$ has a factor $p$.
In the CA chart JD16 divided by $\rho$ has bounded coefficients
on the same raw roots as TL8.  In the SE chart
\[
 \Theta/\rho=\varepsilon^2\tau^2QE/K
\]
gives the same extra $\varepsilon^2$ factor to every new coefficient.
Modulo $\varepsilon$, scalar coefficients are even in $u,v$ and
the normal coefficients are independent of $\mathfrak w$.
Thus the exact support, full-jet, and finite divisibility hypotheses
of CC3--CC6 and SE5--SE8 are unchanged.  The complete stress divided
by $\rho$ has a bounded Gram on the same eight roots.  Multiplication
by $\rho$ makes it an infinitesimal perturbation of the strict formal
C1 base.  Restoring the positive packets proves the actual C2 parent
on the entire thin contact chart, including both full state normals.

At literal $n=0$, OP4 is a physical C1 identification with
$k_1,\zeta_1,r_1$ infinitesimal, $\eta_1=t$ a unit and
$e_{1,\mathrm{scale}}^2=tb$ bounded.  The physical C1.2/C1.3
all-tests theorems allow arbitrary positive $J_1$ and literal beta
zero, so they cover this boundary exactly.

\paragraph{The unit-$W$ entry on the unbounded-$a$ thin layer.}
The actual physical C2 critical coefficient is
\[
 C_*={(h+2Ah)(1+2D)\over2(4AhD-h)}.
\]
Use the polynomial UW3 numerator/denominator with $b=wz_*$,
$z_*=b/w$ bounded.  Its numerator is divisible by $sw$ because it
vanishes identically at $s=0$ and $w=0$.  The denominator after
canceling the $m$ numerator has residue $16t_0=16$, also when
$j_0=0$.  Consequently
\[
                   (C_*-C)/m=s\times O(1).            \tag{JD18}
\]
There is no $j$ inverse in this estimate.  SEO12 permits both $A,D$
unbounded: its entire hard block remains strict and its two radial
masses tend to $2,2$.  With $W=1$, the exact testing identity
$vZ=Z(R_m+X+uZ)/(1+uY)$ absorbs the full negative critical correction
by JD18 and the retained original $m$ square.  All denominators are
pure testing units.  This proves every local normalized unit-$W$
state, including $Z=0$; negative-core state normalization leaves
exactly the $Z=1$, $u,v,\mathfrak w>0$ infinitesimal chart proved
above.  This is the same actual-C2 state entry as UW, with its
source bound checked at the new $t_0=1$ boundary.

\paragraph{Finite remote tests when $e$ is infinitesimal.}
For bounded $a$, $C_*-C$ is infinitesimal because $C_*,C\sim1/2$.
For unbounded $a$ this follows also from JD18 and bounded $m$.
The full critical-wall identities SEO17, SEO18, SEO20, in the
normalized SEO12 frame, are strict whenever at least one finite
angle has nonzero residue.  They explicitly allow $A,D$ unbounded.
The entire correction $-4(C_*-C)_+D_v(uW-Y)^2$ has infinitesimal
bounded entries in that frame.  Retain original $n,m,\alpha$ and
any positive excess.  This proves every finite remote test, with
all states, for the entire $e$-infinitesimal R8 domain.

\paragraph{Finite remote tests when $e$ is a unit, including small $r$.}
Here $a$ is bounded, $m,C$ bounded, $C$ a positive unit, $A,D$
unbounded, and $N_*=n+1/2$ is either a unit or unbounded.  All
normalizations below are pure source.  Let
$\widehat X=\sqrt{N_*}X$, $\widehat Y=\sqrt A Y$, and
$\widehat D=\sqrt D\,R_{D,\mathrm{event}}$.

If $u$ is a bounded unit and $v$ bounded, use the exact inverse
\[
 W={D^{-1/2}\widehat D+A^{-1/2}\widehat Y+vZ
                         -uvN_*^{-1/2}\widehat X\over u}.
                                                               \tag{JD19}
\]
All full coefficients are bounded.  The $\widehat Y,\widehat D$
residue pivots are $(1+U)(1+V),1$, with zero cross residue, where
$U=u_0^2>0,V=v_0^2\ge0$.  For unit $r$, the remaining two-state
matrix before $X$ normalization is the complete original parent at
$Y=R_{D,\mathrm{event}}=0$, namely
\[
 \begin{split}
 &(nD_uD_v+mD_v^2+CD_vUV+D_v/2)X^2\\
 &+2\{-mD_v(V-U)/u-CD_vuV+uD_v/2\}XZ\\
 &+\{m(V-U)^2/U+CD_vV+(U-V)/2\}Z^2.
 \end{split}                                                   \tag{JD20}
\]
Insert the residue source values JD7, with $e=1-r$, $g=1-ae$.
Its determinant has the exact positive expression
\[
 \begin{split}
 U\det={D_uD_v(1+ar)\over4r(1+a)^2g}
 \{e(U-aV)^2+r(g+ar)U^2V+g(g+ar)UV^2\}.             
 \end{split}\tag{JD21}
\]
It is positive for every $U>0,V\ge0$.  The $X$ diagonal is
positive.  Division by the positive unit $N_*$ proves strictness
after normalization.  When $r$ is infinitesimal, divide JD21 by
$N_*\sim n$ before taking residue.  The result is
\[
 U\det_{\rm normalized}
 ={D_uD_v\over2(1+a_0)(1-a_0)}
       \{(U-a_0V)^2+(1-a_0)^2UV^2\}>0,             \tag{JD22}
\]
with $0\le a_0<1$.  The normalized $X$ diagonal tends to $D_uD_v$.
Thus JD19 proves this remote chart for arbitrary relative $n,A,D$
rates, without a comparison between any of them.

It remains to take $u$ infinitesimal and $v$ a bounded unit.  If
$a$ is a unit, $m$ is a unit.  Solve instead for $Z$ by the pure
testing unit $v$.  The full residue is
\[
 m_0(v_0W-N_{*,0}^{-1/2}\widehat X)^2
       +D_v\widehat X^2+D_v\widehat Y^2+\widehat D^2,
                                                               \tag{JD23}
\]
where $N_*^{-1/2}$ has residue zero if $n$ is unbounded.  This is
strict, with the two-state determinant $m_0V D_v>0$.
If $a$ is infinitesimal, use the extra legal root $S=aW$ and put
$R=uW$.  The exact D event gives
$R=vZ-uvX+Y+D^{-1/2}\widehat D$.  Substitute this expression
in every occurrence except the $m$ event, where use
$mR_m^2=(m/a^2)(vS-aX+auvY-auZ)^2$.
All coefficients are bounded, also when $n$ is unbounded after
the displayed $X$ normalization.  Its entire residue is diagonal
on $(S,\widehat X,\widehat Y,\widehat D,Z)$ with pivots
\[
                    ((m/a^2)_0V,D_v,D_v,1,V^2/2)>0.   \tag{JD24}
\]
These identities retain every original event and the complete signed
core; introducing a dependent extra root is harmless.  The finite
separated lift around the real angle residues proves each chart.

\paragraph{Projective endpoints, signs, and the complete join.}
On all of R8, $C$ is a positive unit and $A,D$ are unbounded.
For infinity charts retain the whole original $m,n,\alpha$ packets,
clip $C,A,D$ at suitable positive real lower bounds, and retain their
positive excesses.  The exact homogeneous infinity identities used
in SEO19 and its accompanying BE2 routing require only these lower
bounds and $d=1$; they do not require $D$ bounded or $m$ bounded
below.  In particular the delicate $(u,1/v)=(0,0)$ corner has the
strict selected residue
\[
 2X^2+4c(uW-Y)^2+(4A_*+2)Y^2+4D_*Z^2
\]
on $(X,Y,uW,Z)$ after retaining $m$.  All leftover $W$ occurrences
carry $u$.  The remaining single and double infinities are exactly
the same homogeneous event identities.  Literal endpoint values are
taken in those original homogeneous formulas; no higher-degree
contact graph is substituted at infinity.

For unit $r$, source $a$ is bounded or unbounded.  The bounded case
has $e$ unit or infinitesimal and is covered locally by JD7 or
JD8--JD10.  The unbounded case forces $e$ infinitesimal and is
covered by the thick/positive-thin/literal-zero split above.  These
are exhaustive pure-source alternatives.  All finite remote tests
use JD18--JD24 or the critical wall, and every infinity uses the
homogeneous join.  Original testing sign gauges, literal-axis event
identities, state homogenization, and the real/imaginary split give
all signs and complex states.  All inverses and radicals are pure
source or pure testing.  No independent source/testing rate is ever
compared.

This proves the entire R8 $r$-unit source subdomain all tests and
the remote complement for every R8 source.  If $r$ is infinitesimal,
$e$ is a unit and $g>0$ forces bounded $a$ with $a_0<1$; precisely
these hypotheses, with every other original intersection unchanged,
are those of the joint-$n,A,D$ local lemma.  That lemma supplies a
strict full-parent Gram with residue $(1,1,1,(m/a^2)_0,C_0)$ on
all remaining local tests, including the axes.  Therefore the two
source alternatives $r$ unit or infinitesimal now exhaust R8 with
an all-tests certificate.  Literal $n=0$ was retained in the unit-$r$
proof; it is impossible in the small-$r$ part because $nr$ is a unit.
The whole ordered R8 shape is closed, with no extra source or testing
rate restriction.  The 20 new focused checks in \path{check_joint_D.json}
corroborate the new exact identities; the ordered-field bounds and
finite-conductor arguments are the proof, not an inference from a
numerical sample or a checker count.
